\section*{Remerciements}
\addcontentsline{toc}{section}{Remerciements}

Nous remercions en premier lieu Monsieur Abassi, notre professeur référent, pour avoir accepté de suivre ce projet et d’en être
le tuteur. Son accompagnement tout au long du semestre, ainsi que ses retours sur l’étude \acrshort{CFD}, notamment sur le choix
du domaine d’étude, la stratégie de maillage, les simplifications pertinentes et la cohérence des hypothèses, nous ont permis de
gagner en rigueur et en efficacité tout en conservant une vision d’ensemble sur l’avancement.\\

Nous remercions Malo Amaranti, chef de projet sur la partie fusée, pour le pilotage, la coordination et le suivi global du
développement du véhicule. Nous adressons aussi nos remerciements à l’ensemble des enseignants et intervenants de
l’\acrshort{IPSA}, dont les retours ont contribué à consolider nos choix techniques.\\

Nous remercions également Monsieur Sellami pour la qualité de ses conseils concernant la partie système lors de nos échanges en
projet de prototypage rapide. \\

Nous remercions Maxence Fulconis, Fab Manager de l’\acrshort{IPSA}, pour l’accès aux moyens de fabrication et pour nous avoir
permis d’usiner nos pièces à la \acrshort{CNC} dans de bonnes conditions. Nous remercions également les anciens d’AéroIPSA pour
la transmission de savoir-faire pratique, ainsi que l’ensemble des membres impliqués dans le projet au sein de l’association, qui
ont contribué aux différentes réalisations.\\

Nous remercions aussi le \acrshort{CNES} et \gls{Planète Sciences} pour leur cadrage et leurs conseils. Les échanges et
recommandations, ainsi que le \acrlong{CDC} associé, ont permis de définir des limites réalistes au projet et de guider les choix
de conception, avec l’appui de retours techniques apportés lors des points de suivi.\\

Enfin, nous remercions nos amis et nos proches pour leur soutien tout au long du semestre.

\newpage
